% ==============================================================================
% APÊNDICE B — MUTANTES EQUIVALENTES DO TETO
% ==============================================================================

\section{Mutantes que restam vivos no teto}
\label{ap:equivalentes}

A suíte do teto mata 367 dos 369 mutantes válidos (338 por falha de asserção e 29 por \textit{timeout} legítimo). O \autoref{qua:equivalentes} registra, para cada um dos dois restantes, o argumento de por que nenhum teste pode distingui-lo do programa original, e o \autoref{qua:revisao2} registra os seis mutantes que uma primeira análise havia classificado como inalcançáveis e que a segunda revisão matou por injeção de dependência. O registro completo, com o censo do Stryker e o histórico das revisões, está no arquivo \texttt{runs/ceiling/EQUIVALENT\_MUTANTS.md} do repositório.

\begin{quadro}[htb]
    \centering
    \ABNTEXfontereduzida
    \begin{minipage}{\textwidth}
        \caption{Mutantes equivalentes do teto}
        \label{qua:equivalentes}
        \centering
        \begin{tabular}{p{0.5cm}p{2.6cm}p{3.8cm}p{7.9cm}}
            \toprule
            \textbf{\#} & \textbf{Local} & \textbf{Mutação} & \textbf{Por que nenhum teste distingue} \\
            \midrule
            1 & \texttt{lexer.ts:12} & \texttt{currentChar = '\textbackslash 0'} $\to$ \texttt{""} & Inicializador morto: o construtor chama \texttt{readChar()} imediatamente, que reatribui o campo em todos os caminhos antes de qualquer leitura. O campo é privado e nunca é lido antes dessa reatribuição, então nem injeção o alcança. \\
            2 & \texttt{lexer.ts:109} & \texttt{peekChar()} devolve \texttt{""} em vez de \texttt{'\textbackslash 0'} no fim da entrada & O valor só é usado em comparações com \texttt{'>'} e \texttt{'='} e nos classificadores \texttt{isLetter}, \texttt{isNumber} e \texttt{isDot}. \texttt{"".charCodeAt(0)} é \texttt{NaN} e \texttt{'\textbackslash 0'.charCodeAt(0)} é 0; ambos falham em todas as comparações de faixa, e nenhum dos dois é \texttt{'.'}, \texttt{'>'} ou \texttt{'='}. O método é privado e seu retorno nunca escapa da classe. \\
            \bottomrule
        \end{tabular}
        \legend{Fonte: Própria do autor (2026).}
    \end{minipage}
\end{quadro}

\begin{quadro}[htb]
    \centering
    \ABNTEXfontereduzida
    \begin{minipage}{\textwidth}
        \caption{Mutantes classificados como inalcançáveis na primeira análise e mortos na segunda revisão}
        \label{qua:revisao2}
        \centering
        \begin{tabular}{p{0.7cm}p{3.0cm}p{4.2cm}p{6.9cm}}
            \toprule
            \textbf{\#} & \textbf{Local} & \textbf{Mutação} & \textbf{Teste que o mata} \\
            \midrule
            3--4 & \texttt{evaluator.ts:47--48} & ramo \texttt{default} do \texttt{switch} de operadores binários e sua mensagem & \texttt{parse} envolvido com \texttt{vi.mock} devolve, uma vez, \texttt{BinaryExpression(Not, 1, 2)}; \texttt{evaluate} deve lançar \texttt{Operator not not implemented}. \\
            5--6 & \texttt{evaluator.ts:77--78} & ramo \texttt{default} do \texttt{switch} de tipos de nó e sua mensagem & \texttt{parse} devolve, uma vez, \texttt{UnaryExpression(Minus, 1)}; \texttt{evaluate} deve lançar \texttt{Invalid AST node}. \\
            7--8 & \texttt{parser.ts:143--144} & bloco do \texttt{if (!this.expectPeekToken(Rparn))} e sua mensagem & \texttt{vi.spyOn} em \texttt{Parser.prototype.expectPeekToken} devolve \texttt{undefined} uma vez; \texttt{parse('(1)')} deve lançar \texttt{Expected ) but got \ldots}. \\
            \bottomrule
        \end{tabular}
        \legend{Fonte: Própria do autor (2026).}
    \end{minipage}
\end{quadro}
